

\chapter{Análise de Custos}
\label{chapter:costs}

\begin{introduction}
Este capítulo apresenta a lista de materiais com os respetivos preços de aquisição e compara o custo total do protótipo desenvolvido com soluções comerciais equivalentes.
\end{introduction}

\section{Lista de Materiais e Preços}
\label{sec:lista-materiais}

A Tabela~\ref{tab:bom} apresenta a lista completa de componentes adquiridos para a construção do protótipo, com as quantidades, preços unitários e totais correspondentes. Todos os preços incluem IVA à taxa legal em vigor e correspondem aos valores de aquisição efetiva registados durante o desenvolvimento do projeto. Os componentes eletrónicos foram adquiridos na loja Mauser; os componentes mecânicos do sistema de translação na loja online Fillment3D; e os acoplamentos flexíveis na Amazon Espanha. Os parafusos e porcas utilizados na montagem estavam disponíveis em laboratório e não foram contabilizados no custo total.

A estrutura mecânica --- base, plataforma superior, colunas de suporte com rótulas e suportes dos motores --- foi fabricada por impressão 3D a partir de filamento PLA, tendo sido consumido menos de um quilograma de material ao longo de todo o projeto.

\begin{table}[h]
\centering
\caption{Lista de materiais (\textit{Bill of Materials}) do protótipo da mesa sísmica.}
\label{tab:bom}
\begin{tabular}{p{6.2cm}ccrr}
\hline
\textbf{Componente} & \textbf{Fornecedor} & \textbf{Qty.} & \textbf{P. unit. (\euro{})} & \textbf{Total (\euro{})} \\
\hline
\multicolumn{5}{l}{\textit{Eletrónica}} \\
\hline
Microcontrolador Joy-IT NodeMCU-ESP32-C~\cite{JoyITESP32}
  & Mauser & 1 & 11,55 & 11,55 \\
Kit CNC Shield + 4$\times$ DRV8825 Joy-IT ARD-CNC-Kit2~\cite{JoyITCNCKit2}
  & Mauser & 1 & 21,45 & 21,45 \\
Motor de passo Joy-IT NEMA17-04~\cite{JoyITNEMA17}
  & Mauser & 4 & 12,00 & 48,00 \\
Fonte de alimentação DC 12\,V / 16,5\,A / 200\,W~\cite{FonteMauser}
  & Mauser & 1 & 14,69 & 14,69 \\
Cabos \textit{jumper} Dupont F--F 300\,mm (40 un.)~\cite{JumperFF}
  & Mauser & 1 & 5,20 & 5,20 \\
Cabos \textit{jumper} Dupont M--M 300\,mm (40 un.)~\cite{JumperMM}
  & Mauser & 1 & 7,16 & 7,16 \\
Cabos \textit{jumper} Dupont M--F 300\,mm (40 un.)~\cite{JumperMF}
  & Mauser & 1 & 5,20 & 5,20 \\
\textit{Breadboard} 830 contactos~\cite{BreadboardMauser}
  & Mauser & 2 & 3,24 & 6,48 \\
\hline
\multicolumn{5}{l}{\textit{Mecânica}} \\
\hline
Varão trapezoidal TR8$\times$8, \textit{lead} 8\,mm (cortado a meio)~\cite{Fillment3DVarao}
  & Fillment3D & 2 & 4,54 & 9,07 \\
Porca T8 em POM, \textit{lead} 8\,mm~\cite{Fillment3DPorca}
  & Fillment3D & 4 & 1,95 & 7,80 \\
Acoplamento flexível motor--fuso (pack de 5, utilizados 4)~\cite{AcopladorAmazon}
  & Amazon ES & 1\,pack & 8,99 & 8,99 \\
\hline
\multicolumn{5}{l}{\textit{Fabricação}} \\
\hline
Filamento PLA 1,75\,mm ($<$1\,kg utilizado)
  & --- & $<$1\,kg & 17,00/kg & $<$17,00 \\
\hline
\textbf{TOTAL} & & & & \textbf{$\approx$152,59} \\
\hline
\multicolumn{5}{l}{\footnotesize Parafusos e porcas M3/M4 disponíveis em laboratório --- não contabilizados.} \\
\multicolumn{5}{l}{\footnotesize Preço do filamento contabilizado pelo máximo (1\,kg); custo real inferior.} \\
\end{tabular}
\end{table}

O custo total do protótipo situa-se em aproximadamente 153\,\euro{}, valor que inclui todos os componentes eletrónicos, mecânicos e de cablagem adquiridos especificamente para este projeto. Este valor está muito abaixo do limite de 500\,\euro{} estabelecido como requisito não funcional (RNF1, Capítulo~\ref{chapter:requirements}), confirmando que o objetivo de baixo custo foi amplamente atingido, com uma margem de mais de 3:1 face ao teto definido. Importa salientar que este custo reflete uma única unidade construída sem qualquer otimização de fornecimento; em produções de maior escala ou com acesso a fornecedores de componentes em volume, o custo unitário poderia ser reduzido de forma significativa.

\section{Comparação com Soluções Comerciais}
\label{sec:comparacao-comercial}

A Tabela~\ref{tab:comparacao-custos} apresenta uma comparação direta entre o protótipo desenvolvido e as soluções identificadas no estado da arte (Capítulo~\ref{chapter:stateofart}), em termos de custo, número de eixos, capacidade de reproduzir sismos reais, utilização de impressão 3D e disponibilidade em regime \textit{open-source}.

\begin{table}[h]
\centering
\caption{Comparação entre o protótipo e soluções de mesas sísmicas existentes.}
\label{tab:comparacao-custos}
\begin{tabular}{lccccc}
\hline
\textbf{Solução} & \textbf{Custo} & \textbf{Eixos} & \textbf{Sismos reais} & \textbf{Imp. 3D} & \textbf{\textit{Open-source}} \\
\hline
Quanser Shake Table II XY~\cite{QuanserShakeTableIIXY} & $>$10\,000\,\euro{} & 2 & Sim & Não & Não \\
Quanser Shake Table III XY~\cite{QuanserShakeTableIII} & $>>$10\,000\,\euro{} & 2 & Sim & Não & Não \\
SARSAR~\cite{Damci2018}                                & 770--1\,800\,USD   & 1 & Sim & Não & Parcial \\
Shakebot~\cite{Shakebot2022}                           & 1\,300\,USD        & 1 & Sim & Parcial & Sim \\
Mesa 2-DOF~\cite{2DOF2025}                             & $\sim$2\,500\,USD* & 2 & Sim & Não & Sim \\
\textbf{Este projeto}                                  & \textbf{$\approx$153\,\euro{}} & \textbf{2} & \textbf{Sim} & \textbf{Sim} & \textbf{Sim} \\
\hline
\multicolumn{6}{l}{\footnotesize * Custo estimado incluindo componente já existente em laboratório, avaliado em 1\,000\,USD.} \\
\end{tabular}
\end{table}

Da análise da tabela ressaltam três aspetos diferenciadores do presente projeto. Primeiro, é a única solução identificada que combina simultaneamente as características de ser bi-axial, de custo reduzido, construída com impressão 3D e totalmente \textit{open-source}. Segundo, o rácio de custo face à solução comercial bi-axial mais próxima --- a Quanser Shake Table II XY~\cite{QuanserShakeTableIIXY} --- é superior a 65:1, ou seja, o presente protótipo custou menos de 1,6\% do valor estimado para a solução comercial equivalente. Terceiro, mesmo comparado com projetos académicos bi-axiais como a mesa 2-DOF~\cite{2DOF2025}, o protótipo desenvolvido apresenta um custo total inferior, sendo simultaneamente o único a recorrer à impressão 3D como tecnologia de fabrico estrutural.

É importante enquadrar esta comparação com as diferenças de especificações técnicas entre as soluções. A Quanser Shake Table II XY é um instrumento de precisão metrológica, com sensores de posição de alta resolução, amplificadores de potência integrados e software de análise avançado, desenvolvido para investigação exigente. O protótipo aqui descrito é uma solução de menor precisão, com controlo em malha aberta e sem realimentação de posição, mas adequada para o objetivo que motivou o seu desenvolvimento: apoiar o trabalho de avaliação de risco sísmico em países em desenvolvimento, onde o acesso a equipamento comercial é economicamente inviável.
